\subsection{规则分类与RIPPER}
\label{sec:classification-rules}
\label{sec:classification-ripper}

一棵树把全部输入组织成从根出发的路径，规则列表则可以直接列出若干足以作出类别判断的条件。例如，目标类若由几个分离的小区域组成，我们可以分别为这些区域学习短规则，而不要求它们共享同一个根测试。规则简化后还可能重叠，因而需要约定命中多条规则时怎样预测。第\ref{sec:classic-overview-tree-rules}节已介绍这种表示差别；这里采用有序规则，按首次命中的结论判类，没有命中时使用默认类。

逐条学习规则时，一个直接困难是：为了覆盖少量例外，不断增加条件会产生庞大的规则集；过早删去条件，又可能丢失有用的局部结构。早期的AQ系列探索了从样本归纳条件规则的路线，后来的规则学习把覆盖策略与误差控制结合起来。Fürnkranz与Widmer在1994年提出\term{增量减错剪枝}（incremental reduced error pruning, IREP），将单条规则的增长和立即剪枝嵌入逐条学习过程，以避免先生成庞大规则集再统一简化的开销\scite{furnkranz1994irep}。

Cohen在1995年改进IREP的剪枝评分和停止准则，并加入整组规则的反复优化，形成RIPPER（Repeated Incremental Pruning to Produce Error Reduction）。RIPPER通过规则增长、留出剪枝、顺序覆盖和整组优化，在长度、覆盖范围与错误之间实施启发式搜索。下面分解这些步骤\scite{cohen1995ripper}。

\subsubsection{模型结构}
\label{sec:classification-ripper-structure}

RIPPER输出有序规则以及一个默认类别$c_0$。第$r$条规则的前件可写为
\begin{equation}
  a_r(\bm{x})=\bigwedge_{j=1}^{m_r}b_{r,j}(\bm{x}),
  \qquad a_r(\bm{x})\Longrightarrow y=c_r,
  \label{eq:classification-rule-form}
\end{equation}
其中$b_{r,j}$是一个返回真或假的测试，$m_r$为该规则的条件数。例如，$x_1>60$与$x_2=\text{是}$可以组成一个合取，但每个测试的业务含义必须由特征定义决定。规则\term{覆盖}（cover）样本，指样本满足其前件，与样本真实标签是否等于后件无关。

若规则数为$s$，预测时寻找最小的$r\in\{1,\ldots,s\}$使$a_r(\bm{x})$成立，输出$c_r$；不存在这样的$r$时输出$c_0$。二分类中，可以为目标类学习若干规则，把另一类作为默认类。多分类时，经典RIPPER通常按类别出现频率由低到高学习，逐次区分当前类与尚待区分的其他类，把最常见类放在默认位置。这种类别排序有助于为少数类显式生成规则，但不构成少数类召回率的保证。

\subsubsection{参数学习}
\label{sec:classification-ripper-learning}

RIPPER采用\term{顺序覆盖}（sequential covering），这个名字描述了规则逐条“占领地盘”的过程：每次只在还没有被任何规则覆盖的样本上学习一条规则，把它加到有序列表的末尾，然后把它命中的样本从工作集中拿走，再对剩下的样本学习下一条；如此反复，直到没有目标类样本剩下。需要强调的是，被一条规则误覆盖的负例虽然从“找下一条规则”的工作集中移除，但在给整组规则打分时仍然计入，因此一条乱覆盖负例的规则不会因为样本被移走就逃避惩罚。设当前目标类别为$c_+$，属于该类的记录称为正例，其余为负例。保留原始二类训练集$D^{(0)}$用于整体评分，另用$U$表示不断缩小的覆盖工作集。从$U$移除的误覆盖负例仍保留在$D^{(0)}$中，并计入整体错误。

学习单条规则时，先分层随机划分增长集与剪枝集，经典设置约为$2:1$。增长从空前件开始，空前件覆盖所有样本；每次添加一个测试，使前件更严格、覆盖范围更小。测试可直接使用连续特征的阈值，无须一律在训练前离散化。Cohen的实现对缺失特征采用保守约定：涉及该特征的测试在缺失记录上不成立，因此规则需要其他已知字段来覆盖这些记录。

增长时使用源自FOIL（First Order Inductive Learner）的信息增益评分。设加入条件前覆盖$p_0$个正例、$q_0$个负例，加入后覆盖$p_1$个正例、$q_1$个负例，其中$q$不表示全章的样本量$n$。在命题规则设定下，评分为
\begin{equation}
  g=p_1\left(
  \log_2\frac{p_1}{p_1+q_1}
  -\log_2\frac{p_0}{p_0+q_0}
  \right).
  \label{eq:classification-foil-gain}
\end{equation}
括号比较加入条件前后的正例纯度，外面的$p_1$避免只看纯度而忽略保留下来的正例数。候选条件必须保留正例并满足覆盖量等实现约束；增长通常持续到不再覆盖增长集负例，或没有可接受的改进条件。噪声样本可能使负例始终被覆盖，此时由可接受条件和覆盖约束结束增长。

剪枝则允许把规则末端的一串条件一起删掉，也就是比较增长过程中得到的各个前缀，而非只尝试删去一个条件。若某个前缀在剪枝集覆盖$p$个正例、$q$个负例，RIPPER的基本评分是
\begin{equation}
  v=\frac{p-q}{p+q},\qquad p+q>0.
  \label{eq:classification-ripper-pruning}
\end{equation}
在分母为正时，它与覆盖样本中的正确率$p/(p+q)$给出相同排序。剪枝选择评分最好的可接受前缀，平局可偏向较短者；覆盖不到剪枝样本时，分数未定义，需执行预先规定的回退或拒绝规则。这一评分改变了原IREP对规则的偏好，尤其避免单凭覆盖大量样本而认可纯度不足的条件。

单条规则看起来合理，也可能使整组规则过于复杂。\term{最小描述长度}（minimum description length, MDL）把这种取舍表达为：用多少比特写下规则，再用多少比特说明规则预测错了哪些训练标签。记有序规则集为$\mathcal{R}$，总描述长度可概括为
\begin{equation}
  L_{\mathrm{total}}(\mathcal{R};D^{(0)})
  =L_{\mathrm{rules}}(\mathcal{R})
   +L_{\mathrm{errors}}(D^{(0)}\mid\mathcal{R}).
  \label{eq:classification-rule-description-length}
\end{equation}
第二项按完整规则集及默认类计算，因而仍包含误覆盖记录。Cohen的实验允许总长度暂时超过此前的最小值，但超过64比特就停止增加规则，再从后向前删除会增加总长度的规则。64比特是该实现允许暂时增加描述长度的预算；规则增加后统一按总长度检查停止和删除。

算法~\ref{alg:classification-ripper}给出二类子问题的主要流程。伪代码中的增长和剪枝采用前述评分，描述长度的编码细节由具体实现确定；多分类时在外层加入类别排序与默认类约定。

\begin{algorithm}[RIPPER的覆盖学习与规则优化]
  \label{alg:classification-ripper}
  \begin{algorithmic}[1]
    \Require 二类样本$D^{(0)}$，目标类$c_+$，优化轮数$b$
    \Ensure 有序规则集$\mathcal{R}$及默认负类
    \State $\mathcal{R}\gets()$
    \Function{Extend}{$\mathcal{R},D^{(0)}$}
      \State $U\gets D^{(0)}$中未被$\mathcal{R}$覆盖的样本
      \State $L_{\min}\gets L_{\mathrm{total}}(\mathcal{R};D^{(0)})$
      \While{$U$中仍有正例}
        \State 将$U$分层随机划为增长集$G$和剪枝集$P$
        \State 从空前件增长规则$a$，按式\eqref{eq:classification-foil-gain}添加条件
        \State 在$P$上按式\eqref{eq:classification-ripper-pruning}选择$a$的前缀
        \If{无可接受规则，或$a$不覆盖$U$中的任何正例}
          \State \textbf{停止本轮覆盖学习}
        \EndIf
        \State 把$a\Longrightarrow c_+$追加到$\mathcal{R}$
        \State $L\gets L_{\mathrm{total}}(\mathcal{R};D^{(0)})$
        \State $L_{\min}\gets\min(L_{\min},L)$
        \If{$L>L_{\min}+64$}
          \State \textbf{停止本轮覆盖学习}
        \EndIf
        \State 从$U$移除被$a$覆盖的全部样本，包括正例和负例
      \EndWhile
      \State 从后向前删除能使总描述长度降低的规则
      \State \Return $\mathcal{R}$
    \EndFunction
    \State $\mathcal{R}\gets\Call{Extend}{\mathcal{R},D^{(0)}}$
    \For{$t=1,\ldots,b$}
      \For{按顺序考察$\mathcal{R}$中的每条规则}
        \State 从空前件生成替代规则，从原前件生成修订规则
        \State 剪枝候选规则时，评价完整规则集在剪枝集上的错误
        \State 用总描述长度在原规则、替代规则、修订规则之间选择
      \EndFor
      \State $\mathcal{R}\gets\Call{Extend}{\mathcal{R},D^{(0)}}$
    \EndFor
    \State \Return $\mathcal{R}$，无规则匹配时预测负类
  \end{algorithmic}
\end{algorithm}

优化阶段是在整组规则的上下文中改写单条规则。替代规则从空前件重新增长，修订规则在原前件上继续添加条件；对候选规则的剪枝不只看它自身的纯度，还看替换后完整有序规则集的错误。随后用MDL选择保留原规则还是采用候选，必要时重新进行规则删除和剩余正例的覆盖学习。原论文把一次优化称为RIPPER，重复两次称为RIPPER2；优化次数规定搜索预算，最终返回已搜索候选中的选定规则集。

\paragraph{搜索能保证到哪里}
\label{sec:classification-ripper-theory}
规则数量和条件长度决定能拼合哪些区域，当前增长、剪枝与替换策略决定实际访问哪些候选。每条有效规则至少覆盖一个尚未移除的正例，再规定有限的条件集合、规则长度和优化轮数，便可保证搜索结束；这不意味着找到了所有规则列表中的最小训练错误或最短描述。MDL与留出剪枝都是当前搜索中的选择依据，最后仍需在独立数据上评价整组有序规则。

\subsubsection{无序规则的组合}

有序规则按首次命中决定输出，排在前面的规则会遮住后面的匹配。Clark与Boswell在1991年对CN2规则学习的改进中研究了无序规则：针对类别分别生成条件，允许多个规则共同参与预测。有序规则以首次命中解决冲突；无序规则则需要组合匹配规则的类别分布或投票权重。无序表示减少了对单一顺序的依赖，却把解释负担转移到重叠和冲突处理上，因此规则学习和预测时的聚合约定需要配套选择\scite{clark1991cn2}。
无序版本改变了重叠区域中的输出组合，不能只按有序列表理解。减少顺序依赖也不意味着预测无需约定：冲突规则怎样合并、没有规则命中时怎样回退，都属于完整的分类器。

\subsubsection{CN2的束搜索}

RIPPER增长一条规则时沿当前条件链继续搜索；若希望同时保留几条条件链，可以比较另一种规则搜索策略。Clark与Niblett在1989年提出的CN2采用这种做法，是另一条直接规则学习路线。它结合AQ的规则覆盖思想与类似ID3的纯度评价，用\term{束搜索}（beam search）寻找条件合取：每轮扩展已有候选，只保留评分最好的$b$个候选继续搜索，其中$b$为束宽。

束宽大于1时，相较于每轮只保留一条条件链的贪心搜索，束搜索可以评价更多候选，也有机会继续探索另一条暂时评分较低的路径。相较于不截断、把全部扩展都保留下来的穷举，限制束宽又能减少后续需要扩展的候选数。这里的“增加计算”与“节省计算”使用了不同的比较基准。

例如，构造一个最多增长3步的搜索，每条候选链在每一步都恰有4个可扩展条件，且不合并重复候选。单链贪心每层评价4个候选，共$4+4+4=12$个；束宽为2时，第一层评价4个并保留2个，后两层各评价8个，共$4+8+8=20$个；不截断时则评价$4+16+64=84$个。这个例子只比较指定搜索树中的候选评价次数，不代表CN2在所有任务上的运行时间比例。

候选截断仍可能丢弃通向最优规则的条件链，因此束搜索不提供全局最优保证。CN2还用显著性筛选限制偶然形成的纯净小区域\scite{clark1989cn2}。

在相同条件全集、前件长度和列表规模下，有序CN2与RIPPER拥有相同的规则表示能力。束宽增加的是可能找到好规则的搜索机会，不新增条件本身无法表达的关系。

\subsection{模型对比与总结}
\label{sec:classification-trees-comparison}
\label{sec:classification-rules-comparison}

表~\ref{tab:classification-tree-comparison}比较局部分裂与整树规模怎样确定。输入包含连续值或缺失字段时，需要检查测试和分支处理；最终条件要供人审阅时，树深、规则数与冲突处理同样影响选择。

ID3、C4.5与CART共享逐层测试的表示，但没有使用同一套搜索准则。C4.5沿用ID3式构造，调整多值特征评分，并补充连续阈值、缺失记录传递与后剪枝。CART把每次分裂统一为二分，再通过代价复杂度路径选择子树。多路测试可以展开成若干二值测试，却可能需要更多节点；因此“二叉”与“多叉”的名称不能直接决定哪个候选树更简洁。输入处理、分裂选择和整树规模都需要在相同任务下比较。

\begin{table}[htbp]
  \centering
  \small
  \begin{tabularx}{\linewidth}{@{}l*{3}{>{\RaggedRight\arraybackslash}X}@{}}
    \toprule
    方法 & 分裂评分 & 基本分支形式 & 规模控制 \\
    \midrule
    ID3 & 信息增益 & 离散取值多路分裂 & 基本版本无后剪枝 \\
    C4.5 & 增益筛选与增益率 & 离散多路、连续阈值二分 & 悲观错误估计剪枝 \\
    CART & 基尼不纯度下降 & 统一二叉分裂 & 代价复杂度路径及验证选择 \\
    \bottomrule
  \end{tabularx}
  \caption{三种决策树在基本分类设定下的比较。输入支持和停止细节随实现版本变化。}
  \label{tab:classification-tree-comparison}
\end{table}

若尚不清楚复杂条件是否必要，可以先建立只依赖一个字段的基线。Holte在1993年研究的\term{单特征规则}（one rule, OneR）为每个字段建立“取值到多数类”的映射，再选择训练错误较少的字段；“one”指一个字段，通常仍有多条取值规则。它便于检查单字段能够解释多少差异，却不能使用不同字段的交互\scite{holte1993oner}。

需要按几个离散字段直接查询类别时，可以采用\term{决策表}（decision table），把所选字段的取值组合作为键，同键记录共用一个预测。Kohavi在1995年研究了用验证误差选择字段子集的办法。查表减少了逐条匹配条件的工作，但字段组合容易稀疏，连续值的分箱和未见组合的回退都需要另行规定\scite{kohavi1995decisiontable}。

当字段取值有限、允许使用全部取值组合且不限制结构规模时，树、规则列表和决策表都能为每个组合指定类别。它们的主要差异在于怎样组织和搜索这些条件：树复用共同的测试前缀，规则可以为不同局部区域采用不同条件顺序，决策表统一使用一组字段作为键。OneR则把全部判断限制在单个字段，不能表达跨字段交互；它适合检验复杂组合是否确有必要。

结构受限或组合未被样本覆盖时，这些表示才会显出不同取舍。树的深度限制会把一些输入留在同一叶内；规则列表需要处理覆盖、顺序与默认类；决策表需要为稀疏和未见键规定回退。应据此比较独立验证误差、条件数量与预测成本，并把条件能否被审阅作为实际约束。
